\runninghead{The Feast and Its Texts}
\section{The Feast of All Saints and Its Texts}
\label{sec:opening}

The Mass of All Saints opens with words of the Church, not of Scripture. Its
Introit antiphon, in the English of the Lasance missal, is ``Let us all rejoice
in the Lord, celebrating a feast in honor of all the saints, in whose solemnity
the angels rejoice and join in praising the Son of God.'' Its verse is the
first of Psalm~32 (33): ``Rejoice in the Lord, O ye just: praise becometh the
upright.'' The Collect says that God has granted us to venerate the merits of
all his saints \latin{sub una celebritate}, in one celebration, and asks that,
\latin{multiplicatis intercessoribus}, with the intercessors multiplied, he
grant the abundance of his mercy that we desire.

The Lesson is John's vision in the seventh chapter of the Apocalypse. An angel
rising from the east with ``the sign of the living God'' tells the four angels
who have power to hurt the earth and the sea to hurt nothing ``till we sign the
servants of our God in their foreheads.'' John hears the number of the sealed,
``an hundred forty-four thousand \dots\ of every tribe of the children of
Israel,'' and hears it again tribe by tribe, twelve thousand from each of
twelve. Then he sees ``a great multitude, which no man could number, of all
nations and tribes and peoples and tongues, standing before the throne and in
sight of the Lamb, clothed with white robes, and palms in their hands,'' crying
``Salvation to our God, who sitteth upon the throne, and to the Lamb.'' The
angels fall on their faces before the throne and give God a sevenfold praise.
The Lesson ends there, at verse~12, before an elder asks John who the multitude
are.

The Gradual is from Psalm~33 (34): ``Fear the Lord, all ye his saints: for
there is no want to them that fear him.'' Of the next verse it sings only the
second half, \latin{Inquirentes autem Dominum, non deficient omni bono}, that
those who seek the Lord shall lack no good; the first half, ``The rich have
wanted, and have suffered hunger,'' is not sung. The Alleluia is the Lord's
invitation in Mt 11:28, ``Come to me all you that labor and are burdened, and I
will refresh you.'' The Gospel is the opening of the Sermon on the Mount
(Mt 5:1--12a). Jesus sees the multitudes, goes up into a mountain, sits, and
teaches: blessed are the poor in spirit, the meek, those who mourn, those who
hunger and thirst after justice, the merciful, the clean of heart, the
peacemakers, and those who suffer persecution for justice' sake; and blessed
are the disciples when they are reviled and persecuted for his sake, ``for
your reward is very great in heaven.''

The Offertory takes three verses of Wisdom and changes them: ``The souls of
the just are in the hand of God, and the torment of malice shall not touch
them: in the sight of the unwise they seemed to die, but they are in peace.
Alleluia'' (Lasance). The Secret offers ``the gifts of our devotion'' and asks
that they be pleasing ``for the honor of all the just'' and profitable to those
who offer them. The Communion sings again three of the beatitudes the Gospel
has just given: the clean of heart, the peacemakers and the persecuted. The
Postcommunion asks that God grant his faithful people ``ever to rejoice in the
veneration of all Thy saints, and to be defended by their unceasing prayers.''

The Missal prints this Mass in the Sanctorale under the heading \latin{Die 1
novembris}, \latin{In festo omnium Sanctorum}, of the first class (printed
pp.~702--703, nos.~3974--3983). In 2026 the first of November is a Sunday, the
Twenty-third after Pentecost. A feast of the first class of the universal
Church outranks a Sunday of the second class (RG~91, rows 11 and 15; RG~16),
and the Sunday is not resumed on another day (RG~14). Because All Saints is not
a feast of the Lord, the Sunday is not simply set aside: it is commemorated, as
a privileged commemoration made in every Mass of the day (RG~108, 109~a,
111~a; RGMR~434~a, 436--437). After the feast's Collect, Secret and
Postcommunion the priest therefore says the Sunday's Collect \latin{Absolve,
quaesumus, Domine}, its Secret \latin{Pro nostrae servitutis augmento} and its
Postcommunion \latin{Quaesumus, omnipotens Deus}, each under its own
conclusion. Nothing else of the Sunday is said. The colour is white
(RG~120~d); the Gloria and the Credo are said (RGMR~431; 475~a and b), and the
last Gospel is the beginning of the Gospel according to St John
(RGMR~509--510). The feast has no Preface of its own. The general rubrics
admit two answers, the Preface of the Most Holy Trinity and the common
Preface, and the grounds of each are given with the appointed texts.

The texts raise three questions, and each of the three interpretations that
follow takes up one. The first is who and how many the saints are. The Lesson
counts the sealed exactly and then shows a multitude that nobody could count.
St Bede takes the number to signify the whole Church, which is beyond
counting; St Augustine, quoting the same multitude, says that the saints, few
beside the lost, are many in the company of the angels, and says elsewhere
that their number is known to the Creator and not to us. The second is why the
Gospel of this feast is the Beatitudes, which bless the poor, the mourning, the
hungry and the persecuted, and in what sense their promises are already given.
St Augustine says that the promises can be fulfilled in this life; St Thomas
Aquinas says the rewards are begun on the way and given perfectly in the
homeland; St Hilary of Poitiers places the vision of God after the change to
immortality. The third is how the Gradual's promise can be true, when saints
have gone hungry and died at persecutors' hands, and what the Offertory means
by the just in the hand of God. St Basil the Great answers that the good they
cannot lack is God himself, and that the goods promised are not bodily ones,
since the apostles sought the Lord and lacked them; St Thomas and St Robert
Bellarmine keep, beside the spiritual sense, a promise of help in bodily
need.

\begin{maptable}
Introit &
\latin{Gaudeamus omnes in Domino}; verse Ps 32:1 (33:1) &
An antiphon of the Church: all rejoice in the Lord on the feast of all the
saints, at whose solemnity the angels rejoice and praise the Son of God; the
psalm bids the just rejoice in the Lord.\\
Collect &
\latin{Omnipotens sempiterne Deus, qui nos omnium Sanctorum tuorum merita} &
God has given us to venerate the merits of all his saints in one
celebration; that, the intercessors being many, he grant the abundance of
his mercy.\\
Commemoration Collect &
\latin{Absolve, quaesumus, Domine}; Twenty-third Sunday after Pentecost &
That God loose the sins of his peoples and free them from the bonds their
frailty has contracted.\\
Lesson &
Apoc 7:2--12 &
The servants of God sealed on their foreheads; a hundred and forty-four
thousand from the twelve tribes; then a multitude no one could number, from
every nation, before the throne and the Lamb; the angels' sevenfold praise.\\
Gradual &
\latin{Timete Dominum, omnes Sancti eius}; Ps 33:10, 11b (34:10, 11b) &
Fear the Lord, all his saints, for nothing is lacking to those who fear him;
those who seek the Lord shall lack no good.\\
Alleluia &
\latin{Venite ad me omnes}; Mt 11:28 &
The Lord calls all who labour and are burdened, and promises to refresh
them.\\
Gospel &
Mt 5:1--12a &
On the mountain Jesus teaches the eight beatitudes, and tells the disciples,
reviled and persecuted for his sake, that their reward is great in heaven.\\
Offertory &
\latin{Iustorum animae in manu Dei sunt}; Wis 3:1--3, adapted &
The souls of the just are in God's hand, and the torment of malice shall not
touch them; to the unwise they seemed to die, but they are in peace;
alleluia.\\
Secret &
\latin{Munera tibi, Domine, nostrae devotionis offerimus} &
The gifts of our devotion, offered that they may be pleasing for the honour
of all the just and, by God's mercy, salutary to us.\\
Commemoration Secret &
\latin{Pro nostrae servitutis augmento}; Twenty-third Sunday after Pentecost &
The sacrifice of praise offered for the increase of our service; that God
complete what he has given to the undeserving.\\
Preface &
The Trinity Preface (no.~1082) or the common Preface (no.~1086) &
The feast has no Preface of its own; the general rubrics can be read for
either.\\
Communion &
\latin{Beati mundo corde}; Mt 5:8--10 &
The clean of heart shall see God; the peacemakers shall be called children of
God; the persecuted for justice possess the kingdom of heaven.\\
Postcommunion &
\latin{Da, quaesumus, Domine, fidelibus populis} &
That the faithful peoples ever rejoice in venerating all the saints and be
defended by their unceasing prayer.\\
Commemoration Postcommunion &
\latin{Quaesumus, omnipotens Deus: ut, quos divina}; Twenty-third Sunday after
Pentecost &
That those whom God gives to rejoice in sharing divine things be not left
subject to human dangers.\\
\end{maptable}
